
\AddSymbolsB
\pagecolor{cyan!10!gray}
\color{white}
\quad\quad \lettrine[,lraise=0.2,lhang=0.5]{\textbf{盛}}{夏}的夜晚，星空璀璨。

伊藤佐一正站在草地上，仰望着天空。

今天，家里没人。没有人陪伴，也没有人呼唤。伊藤佐一感受到了一种疼痛，从他的脑中蔓延，接着便长满他的头皮。

\twkkai{应该是风吹的吧；}他闭上了眼睛，开始感受起风的行迹。

公园里只有寥寥几个人，可是马路上却热闹极了——车流涌动，似乎永不会停息。小鸟依旧被街灯高高挂起，而街灯仿佛为鸟儿准备了一个舞台，向花丛打着光。

鸟叫了几声。伊藤佐一听着。

而他之所以仍有心情聆听，只是因为铃木和子今天找过他。她只说了短短一句话。而他仍记得。

他回忆着。

\twkkai{如果可以，你以后能帮我抄抄笔记吗？}

她说这句话的时候，走得是那么近。

\twkkai{为什么是我呢？}他想着。他不曾怀疑他内心的舒畅；他也从未质疑过她带给他的“麻烦”。

\twkkai{难道这是上天安排的羁绊吗？}他沉思着。

远处，公园的长椅似乎在诱惑着他，促使他的双腿疲软。他走向了长椅，打算坐下——以便继续思考。

走着走着，他便迎头撞上一个抱着孩子的年轻母亲。

\twkkai{啊！不好意思，不好意思……}伊藤佐一捂着他的头，低头猛地道歉。

\twkkai{没事……我想问问，地铁站怎么走？}那位母亲问道。

伊藤佐一指了指方向。指的时候，他突然想起那座亲切的图书馆，而忘记了那条长椅。

\twkkai{是啊，为什么不去那里呢——那一切真正发生的地方……；}伊藤佐一发着呆，想着。

女人回头看向远方的地铁口，便向伊藤佐一道了谢，离开了。等女人走远后，伊藤佐一也走向了地铁口。

伊藤佐一来到地下平台。他看着地铁玻璃门里的自己。那个自己，被飞驰而过的列车穿透、打破，变得支离破碎；转而变成的，是众生的脸。那是一张张受了奔波的伤害而痛苦的脸，它们低垂着，仿佛随时要融化。

伊藤佐一早已适应了这一切——每天都是这样；只有铃木和子可以让这一切变得新奇一点。

地铁门开了。他走进了车厢，找了最近的一个空座位，坐了下来。车厢门履行着它的职责，缓缓地关上了。车厢摇晃着前进；窗外闪过一幅幅炫目的广告牌，刊登着不同的产品、人物。伊藤佐一早已见过每一幅广告牌，对它们并不感兴趣。

他能逐渐感受到一股向上升的加速度。列车离开了地底隧道，来到了城市上方的轨道桥上。摩天楼拔地而起；霓虹灯的光芒在空气中散射着；汽车呼啸而过。

伊藤佐一觉得这一切都是那么不真实，可是此时的他被迫接受着周遭世界的真实性。

那种真实性逼他走到了图书馆。在他意料之外的是，图书馆仍亮着灯，而且还有很多人在埋头自习。

伊藤佐一走回了那个一切惊喜发生的书架。没有人。

他回到之前那个熟悉的位置，从标满记号的书架上拿出一本书。那是教授推荐看的书。

他翻开书来，却掉落了一张卡纸。于是他便蹲在地上，并不把它捡起，只是静静地看着。

可是蹲到大腿发麻了，他还是没看出个所以然来。

\newpage
\color{black}
\pagecolor{white}
\twkkai{\textbf{临时草稿，暂时保密：}
\begin{thrm}[Landau\twkkai{定理}]\label{landau}\index{landau@Landau 定理}
\twkkai{
设 $f(s) = \displaystyle\sum_{n=1}^{\infty} \dfrac{a_n}{n^s}$ 是系数 $a_n \geq 0$ 的狄利克雷级数，收敛半平面为 $\text{Re}(s) > \sigma_0$。则 $s = \sigma_0$（实轴上的边界点）必是 $f(s)$ 的奇点。}\index{dilikeleijishu@狄利克雷级数}\end{thrm}

\begin{proof}
\twkkai{
不妨设 $\sigma_0 = 0$（用 $s - \sigma_0$ 替换 $s$ 即可），即级数在 $\text{Re}(s) > 0$ 收敛，且 $f$ 能解析延拓至 $s=0$ 的邻域。

由假设，$f(s)$ 在 $\text{Re}(s) > 0$ 及 $s = 0$ 的邻域内均解析，故存在 $\varepsilon > 0$，使得 $f(s)$ 在圆盘
$$D = \{s \in \mathbb{C} : |s - 1| \leq 1 + \varepsilon\}$$
内解析（注意 $s = -\varepsilon$ 满足 $|-\varepsilon - 1| = 1+\varepsilon$，恰在圆盘边界上）。

$f$ 在 $\text{Re}(s) > 0$ 内可对级数逐项求导：
$$f^{(p)}(s) = \sum_{n=1}^{\infty} \frac{a_n (-\log n)^p}{n^s}, \quad \text{Re}(s) > 0.$$

令 $s = 1$：
$$f^{(p)}(1) = (-1)^p \sum_{n=1}^{\infty} \frac{a_n (\log n)^p}{n}.$$
因 $a_n \geq 0$，故 $(-1)^p f^{(p)}(1) \geq 0$。

$f$ 在 $D$ 内的 Taylor 展开为：
$$f(s) = \sum_{p=0}^{\infty} \frac{f^{(p)}(1)}{p!}(s-1)^p, \quad |s-1| \leq 1+\varepsilon.$$

代入 $s = -\varepsilon$后，由于 $|-\varepsilon - 1| = 1 + \varepsilon$，点 $s = -\varepsilon$ 在 $D$ 的边界上，Taylor 级数仍收敛：
$$f(-\varepsilon) = \sum_{p=0}^{\infty} \frac{f^{(p)}(1)}{p!}(-1-\varepsilon)^p = \sum_{p=0}^{\infty} \frac{(-1)^p f^{(p)}(1)}{p!}(1+\varepsilon)^p.$$

因各项 $\dfrac{(-1)^p f^{(p)}(1)}{p!}(1+\varepsilon)^p \geq 0$，由 Tonelli 定理可交换求和顺序：
$$f(-\varepsilon) = \sum_{p=0}^{\infty} \sum_{n=1}^{\infty} \frac{a_n (\log n)^p (1+\varepsilon)^p}{p! \cdot n} = \sum_{n=1}^{\infty} \frac{a_n}{n} \sum_{p=0}^{\infty} \frac{\bigl((1+\varepsilon)\log n\bigr)^p}{p!}.$$

注意到内层求和恰为指数函数的 Taylor 级数：
$$\sum_{p=0}^{\infty} \frac{\bigl((1+\varepsilon)\log n\bigr)^p}{p!} = e^{(1+\varepsilon)\log n} = n^{1+\varepsilon}.$$

故：
$$f(-\varepsilon) = \sum_{n=1}^{\infty} \frac{a_n}{n} \cdot n^{1+\varepsilon} = \sum_{n=1}^{\infty} \frac{a_n}{n^{-\varepsilon}}.$$

结论：上式表明狄利克雷级数 $\sum a_n / n^s$ 在 $s = -\varepsilon$ 处收敛。  
由狄利克雷级数收敛半平面定理，级数在整个半平面 $\text{Re}(s) > -\varepsilon$ 内收敛。
这与"$\sigma_0 = 0$ 是收敛横坐标"矛盾，故 $s = 0$ 必为 $f(s)$ 的奇点。}
\end{proof}

\begin{rmk}
\twkkai{条件 $a_n \geq 0$ 是本质的，它保证了交换求和顺序的合法性，也保证了 $(-1)^p f^{(p)}(1) \geq 0$。若去掉此条件，定理一般不成立（如黎曼$\zeta$ 函数的狄利克雷展开可解析延拓越过收敛边界）。}\end{rmk}



}
\clearpage
\pagecolor{cyan!10!gray}
\color{white}

\twkkai{看着那么像教授的字迹——是他写的吗？}

伊藤佐一捡起了那张卡纸，阅读了起来。他想起了铃木和子的话。

\twkkai{不如趁现在，把它抄下来吧；}伊藤佐一想着。可是他才猛然看见“\twkkai{临时草稿，暂时保密}”八个大字。

\twkkai{不管了——答应了就要去做；}他还是选择笃定自己的决心。

他依旧悄悄地踮起脚，走进自习的人群中，找到一个空座位坐下，慢慢誊写起来。

谁知，抄了没一会儿，伊藤佐一仿佛像后脑勺受到了重击一般，晕倒在了地上。

……

铃木和子正在操场上散着步。

她低着头，为今天所发生的事闷闷不乐。实际上，今天什么都没有发生——除了伊藤佐一无故的缺席。

她走上课室阶梯的时候，被人推搡着。这番情景让她猛然觉得——平日里在她眼中略显毛手毛脚的伊藤佐一，若放在人群中对比，一定会变得过于拘谨。

今天早上，回到了课室的最后一排后，铃木和子安静地坐了下来，拿出了那张空白的试卷；她立刻觉得烦闷起来——空气是那么地潮湿而令人窒息。

上课17分钟了；可是属于伊藤佐一的座位依然空着。而铃木和子的心此时也悬着。她不确定——她不敢确信自己的期待，也不敢期待伊藤佐一的到来；毕竟，她已经两个星期没头痛了——当然是不想再有第三次的了。

风打扰了铃木和子的思绪，吹起了铃木和子的刘海；她猛地把刘海从头顶轻轻地挡下来。

突然，操场上传来一阵欢呼声。铃木和子用手遮住刘海，抬头望着——天空中飘来几十朵奇形怪状的云，让她惊异无比：

\noindent\adjustbox{center}{\includegraphics[scale=1.3]{plot.pdf}}

铃木和子把头低了下来，又不安地看向人群。

这时，有人从身后拍了拍她的肩膀。铃木和子吓得扭回了头，眼望前方，不敢回头。

她期待是他，又祈祷着不是他。

\twkkai{铃木同学？}

在她耳畔传来的是一个熟悉的女孩的声音。她回过头来，才略显失望地发现，原来是社团的那位女同学；但她显然不敢将她的失望明目张胆地摆在脸上。

\twkkai{怎么啦？}

\twkkai{这是你的吗？刚才有一只猫一直跟着你，在你脚后跟放东西呢；}那位女同学说着，递出了一块奖章。

铃木和子立刻警惕起来；但她的不好意思让她伸出了手。

\twkkai{谢谢你；}她不敢相信，自己竟让这一切发生。她的手沉了一下。

\twkkai{记得保管好哦，先走了……哦对了！明天有社团活动，记得过来参加！}女同学把奖章轻轻放进铃木和子的手里，便和她告别，跟远处的人群汇合了。

那几十朵细碎的云还没有消散，仍悠闲而带着威严地悬挂在天空。

\begin{center}
——————
\end{center}

教授拖着伊藤佐一，在一座教堂里漫步着。教堂的一处窗沿站着几只傻乎乎地探着头的鸽子。这当然让刚睡醒的伊藤佐一疑惑不已。

\twkkai{教授？}他小声喊着；但教授显然没听见。

\twkkai{教授！}他转而大声喊着；他不曾想，自己竟敢对教授如此粗鲁。他能感受到自己停了下来——停得是那么急促，仿佛此时此刻是他反向拖着教授向前走。

\twkkai{我们这是去哪？我为什么在这里？你为什么这样拖着我？}伊藤佐一急促地问着，生怕教授下一刻就把他扔下不管，逃之夭夭。

教授自然没有回答。伊藤佐一开始急速思考。

\twkkai{这一切都太不合理了——教授不可能这么做，我也不可能这么说；我们更不会同时出现在教堂……教堂……}

他扭头观察着。四周的混凝土墙镶嵌着十几扇彩色的琉璃玻璃窗，中间的十字支撑架把窗上的光搅成四瓣；眼前是一条极长的走廊，两旁的长椅——仿佛就像公园的那些长椅——并没有坐着任何一个忏悔的人。

但他瞬间看到椅子下多了许多铁制十字架。

\twkkai{这是一场梦吧，何不……？}

他立即抄起其中一条十字架，砸向自己的头颅。无事发生；但他依旧拿着那沉甸甸的十字架。

教授依旧头也不回地把他往教堂讲台的方向拖着。伊藤佐一再次重击自己的头颅。

\twkkai{停止你的所作所为！}教授厉声喝道。窗沿的鸽子吓得四处逃散，占满了整片天空。

……

铃木和子惊醒了。操场上的人儿依旧热闹地在那几十朵细碎的云下奔跑。

她才意识到，自己经历了最近一个月以来第三次最剧烈的头痛。

\twkkai{原来头痛是这么缓解的；}她摇了摇头，确认自己不再头痛，又同时表示着自己的无奈。

\twkkai{那头痛一定是来自……；}她看了看手里的奖章。





















